\documentclass[10pt,a4paper]{amsart}
\usepackage[margin=19mm]{geometry}
\usepackage{amsmath,amssymb,mathtools}
\usepackage[scr=rsfs,cal=euler]{mathalfa}
\usepackage{xfrac}
\usepackage[unicode,hidelinks,bookmarks=false]{hyperref}
\newcommand{\ie}{i.e.\ }
\newcommand{\cf}{cf.\ }
\newcommand{\resp}{respectively\ }
\newcommand{\Subsection}{\subsection}
\providecommand{\nobreakdash}{\nobreak\hbox{-}}
\setlength{\emergencystretch}{2em}
\multlinegap0pt
\usepackage{bm}
\usepackage{bbm}
\usepackage{dsfont}
\usepackage{xspace}
\usepackage{cancel}
\usepackage{mathtools}
\usepackage[shortlabels]{enumitem}
\usepackage{amsthm}
\makeatletter
\@ifundefined{theorem}{\newtheorem{theorem}{Theorem}}{}
\@ifundefined{proposition}{\newtheorem{proposition}[theorem]{Proposition}}{}
\makeatother
\title[The Extended Real Line with Reentry: Separating US from KC i]{The Extended Real Line with Reentry: Separating US from KC in the Clontz Hierarchy}
\author{Damian Rafael Lattenero}
\date{Source: arXiv:2603.03228v5 (CC BY 4.0)}
\begin{document}
\maketitle

\noindent\emph{Source and licence.} The Extended Real Line with Reentry: Separating US from KC in the Clontz Hierarchy. Version arXiv:2603.03228v5 (CC BY 4.0), distributed under the Creative Commons Attribution 4.0 International licence, as recorded in the arXiv licence metadata for this exact version and shown on the abstract page https://arxiv.org/abs/2603.03228v5. Full rights notice: licenses/arxiv-2603.03228-CC-BY-4.0.txt.\par\smallskip

\noindent\textit{Editorial scope.} Counted unit: the Theorem whose source label is \texttt{thm:sober} in the article (source0.tex lines 1247--1249, source label \texttt{thm:sober}), delivered together with the article's own complete original proof of it (source0.tex lines 1251--1346, one \texttt{proof} environment of 96 lines). No mathematics is written, repaired or re-derived: statement and proof are the article's own text line for line. Required context retained uncounted: prop:T1 (lines 364--366). Every definition, display and label of those lines is retained unchanged. Omitted from this package and not redistributed: 1--363, 367--1246, 1250--1250, 1347--3255. Nothing inside a retained excerpt is omitted or altered: every retained body is delivered exactly as the article's lines are written, so no omission receipt of any kind accompanies this package. Citations: the retained excerpts carry no citation command at all, so no bibliography entry of the article is transcribed and no reference is redistributed. Reference adaptation: none was needed and none was made; every reference inside the delivered unit resolves inside it, so no unresolved reference or citation is produced. Printed slips kept literal: no printed word, symbol, hypothesis or number of the counted statement or of its proof was altered. Layout and class: the article's own class is \texttt{article} with packages including inputenc, fontenc, babel, amsthm, amsmath, amssymb, geometry, enumitem, microtype, xcolor, url, hyperref, tikz; the wrapper is the lane's pinned \texttt{amsart} editorial wrapper at 10pt on A4, which restates the theorem-like environments the retained text needs and adds the article's own macro definitions that the retained text needs (none), with extra wrapper packages (bm, bbm, dsfont, xspace, cancel, mathtools, enumitem). The delivered numbering is the unit's own. Assets: no retained span contains a figure environment, a graphics-inclusion command, a PDF-page inclusion, a table or a TikZ picture; no image file is copied into this package, the package declares no asset dependency and no image file is redistributed with it. Licence: version arXiv:2603.03228v5 of this article is distributed under the Creative Commons Attribution 4.0 International licence, as recorded in the arXiv licence metadata for this exact version and shown on the abstract page https://arxiv.org/abs/2603.03228v5. A full rights notice is delivered at licenses/arxiv-2603.03228-CC-BY-4.0.txt. Characters: every retained excerpt is pure ASCII, so the delivered engine, XeLaTeX with its default Latin Modern text font, reports no missing character.
\begin{proposition}[$X$ is $T_1$]\label{prop:T1}
Every singleton in $X$ is closed.
\end{proposition}

\begin{theorem}[Sobriety]\label{thm:sober}
$(X, \tau_{\mathrm{ERI}})$ is sober.
\end{theorem}

\begin{proof}
A space is \emph{sober} if every irreducible closed subset has a unique
generic point.  A closed set~$F$ is \emph{irreducible} if whenever
$F = F_1 \cup F_2$ with $F_1, F_2$ closed, then $F_1 = F$ or $F_2 = F$.
A point~$x$ is \emph{generic} for~$F$ if $\overline{\{x\}} = F$.

Since $X$ is~$T_1$ (Proposition~\ref{prop:T1}), every singleton is
closed, so $\overline{\{x\}} = \{x\}$ for all $x \in X$.  Hence a closed
set can have a generic point only if it is a singleton, and each
singleton~$\{x\}$ has the unique generic point~$x$.  Sobriety therefore
reduces to the following claim.

\medskip
\noindent\textbf{Claim.}\; \emph{Every closed subset $F \subseteq X$
with $|F| \geq 2$ is reducible} (i.e., expressible as a union of two
proper closed subsets).
\medskip

\noindent\emph{Case 1: $\ast \notin F$.}\;
Then $F \subseteq X \setminus \{\ast\}$.  The subspace
$X \setminus \{\ast\}$ carries the Euclidean topology of
$\mathbb{R}\setminus\{0\}$ (since $q$ restricts to a homeomorphism
$\mathbb{R}\setminus\{0\} \xrightarrow{\;\sim\;} X\setminus\{\ast\}$),
hence is Hausdorff.  Every Hausdorff space is sober: if $|F| \geq 2$,
pick distinct $a, b \in F$ and separate them by disjoint open sets
$U_a, U_b$ in the subspace; then $F = (F \setminus U_a) \cup (F \setminus U_b)$
is a decomposition into two proper closed subsets of~$F$.
Since irreducibility is intrinsic to the subspace topology on~$F$,
the set~$F$ is reducible.

\smallskip
\noindent\emph{Case 2: $\ast \in F$ and $|F| \geq 2$.}\;
Pick $\hat{c} \in F \setminus \{\ast\}$ with
$c \in \mathbb{R}\setminus\{0\}$.  Choose $\varepsilon > 0$ small
enough that $[c - \varepsilon,\, c + \varepsilon] \subseteq
\mathbb{R}\setminus\{0\}$ (possible since $c \neq 0$).  Define
\[
  K \;:=\; q\bigl([c-\varepsilon,\, c+\varepsilon]\bigr),
  \qquad
  L \;:=\; X \setminus q\bigl((c-\varepsilon,\, c+\varepsilon)\bigr).
\]
We verify that $K$ and $L$ are closed in~$X$.

\emph{$K$ is closed.}\;
Its complement satisfies
$q^{-1}(X \setminus K) = \overline{\mathbb{R}} \setminus
[c-\varepsilon, c+\varepsilon]$, which is open in~$\overline{\mathbb{R}}$.
Moreover, $\{-\infty, 0, +\infty\} \subseteq
\overline{\mathbb{R}} \setminus [c-\varepsilon,c+\varepsilon]$
(since $[c-\varepsilon,c+\varepsilon] \subseteq
\mathbb{R}\setminus\{0\}$), and
$\overline{\mathbb{R}} \setminus [c-\varepsilon,c+\varepsilon]$
is dense in~$\overline{\mathbb{R}}$ because it contains
$(-\infty, c-\varepsilon) \cup (c+\varepsilon, +\infty)$.
Since $\ast \in X \setminus K$, condition~(b) requires this density,
which holds.  So $X \setminus K \in \tau_{\mathrm{ERI}}$.

\emph{$L$ is closed.}\;
Its complement satisfies
$q^{-1}(X \setminus L) = (c-\varepsilon, c+\varepsilon)$, which is open
in~$\overline{\mathbb{R}}$.
Since $\ast \notin X \setminus L$ (as
$\{-\infty, 0, +\infty\} \cap (c-\varepsilon,c+\varepsilon) = \varnothing$),
condition~(b) is vacuous.  So $X \setminus L \in \tau_{\mathrm{ERI}}$.

Now set $F_1 := F \cap K$ and $F_2 := F \cap L$.  Both are closed
(intersections of closed sets).  We check the three required properties.
\begin{itemize}[nosep]
\item \emph{$F_1 \subsetneq F$:}\;
  Since $q^{-1}(K) = [c-\varepsilon,c+\varepsilon]$ does not meet
  $\{-\infty,0,+\infty\}$, we have $\ast \notin K$, hence
  $\ast \notin F_1$.  But $\ast \in F$, so $F_1 \subsetneq F$.
  Also $\hat{c} \in K \cap F$ gives $F_1 \neq \varnothing$.
\item \emph{$F_2 \subsetneq F$:}\;
  Since $c \in (c-\varepsilon,c+\varepsilon)$, we have
  $\hat{c} \in q\bigl((c-\varepsilon,c+\varepsilon)\bigr)$, hence
  $\hat{c} \notin L$ and $\hat{c} \notin F_2$.  But $\hat{c} \in F$,
  so $F_2 \subsetneq F$.
  Also $\ast \notin q\bigl((c-\varepsilon,c+\varepsilon)\bigr)$ gives
  $\ast \in L$, so $\ast \in F_2 \neq \varnothing$.
\item \emph{$F_1 \cup F_2 = F$:}\;
  It suffices to show $K \cup L = X$.  Let $\hat{x} \in X$.  If
  $\hat{x} \in K$ we are done.  If $\hat{x} \notin K$, then
  $x \notin [c-\varepsilon, c+\varepsilon]$, so in particular
  $x \notin (c-\varepsilon, c+\varepsilon)$, hence
  $\hat{x} \notin q\bigl((c-\varepsilon,c+\varepsilon)\bigr)$, giving
  $\hat{x} \in L$.
\end{itemize}
Thus $F = F_1 \cup F_2$ with both $F_i$ proper closed subsets
of~$F$; hence $F$ is reducible.

\medskip
In both cases every closed set with $|F| \geq 2$ is reducible.
The only irreducible closed sets are the singletons~$\{x\}$, each with
unique generic point~$x$.  Hence $(X, \tau_{\mathrm{ERI}})$ is sober.
\end{proof}

\end{document}
